\section{Related Work}
\label{sec:related}

We position our work along three axes: privacy attacks against FL,
privacy mechanisms for FL, and the largely separate literature on
topology in FL. The gap our work fills, namely topology-conditional
distributional inference under DP with a defense paired by construction,
sits at the intersection of all three and, to our knowledge, has
not been addressed.

\subsection{Privacy Attacks Against Federated Learning}
\label{sec:related-attacks}

Two attack families dominate the literature, both content-oriented.
\emph{Gradient inversion} reconstructs training samples by optimising
dummy inputs to match observed gradients. Zhu et~al.\ introduced Deep
Leakage from Gradients (DLG)~\cite{zhu2019dlg}, which recovers
pixel-accurate images from per-batch gradients in the absence of noise;
Geiping et~al.\ subsequently demonstrated that magnitude-invariant
formulations can recover images even from federated-averaging updates
covering multiple local steps~\cite{geiping2020inverting}. The fundamental
constraint of this family is sensitivity to per-batch noise: in the
DP-SGD regime that we study, gradient inversion degrades sharply because
clipping plus Gaussian perturbation removes the high-frequency signal
needed for pixel-level reconstruction. We treat gradient inversion as a
\emph{complementary} threat that is most relevant when DP is weak;
under the DP-SGD regime we study, it is outside the scope of our
empirical evaluation.

\emph{Inference attacks} target distributional or membership properties
rather than per-sample reconstruction. Shokri et~al.\ established the
shadow-model paradigm for membership inference~\cite{shokri2017membership},
in which the adversary trains an attack model on shadow datasets where
membership is known and applies it to the target. Melis et~al.\ extended
this to the collaborative-learning setting and showed that joint update
sequences leak unintended properties about subsets of training
data~\cite{melis2019exploiting}. Our \TADI attack inherits the
shadow-training framework but targets a per-client, distributional
quantity (sensitive-class concentration $p_i$) and explicitly decomposes
its information channels into parameter, structural-position, and
organisational-label components. To our knowledge no prior attack
isolates these channels.

\subsection{Privacy Mechanisms for Federated Learning}
\label{sec:related-defense}

The dominant defense is differential privacy applied at the client.
DP-SGD~\cite{abadi2016dpsgd} clips per-record gradients and adds
calibrated Gaussian noise; McMahan et~al.\ instantiated DP-SGD inside
federated averaging as DP-FedAvg~\cite{mcmahan2018dpfedavg} and
demonstrated its viability for production-scale next-word prediction.
Wei et~al.\ provided the first comprehensive convergence analysis of
DP-FL via the noising-before-aggregation framework~\cite{wei2020nbafl},
which we adapt for our utility analysis (Proposition~\ref{prop:utility}).
Privacy accounting in this space has matured around R\'enyi
DP~\cite{mironov2017rdp}, with subsampling
amplification~\cite{wang2019subsampled} extending the formalism to the
Poisson-batched setting that DP-SGD requires. The mutual-information
characterisation of DP~\cite{cuff2016mi} and its tighter conversions
between RDP and other variants~\cite{asoodeh2021variants} provide the
machinery we use to convert the per-mechanism RDP bound into the
per-client MI bound of Theorem~\ref{thm:per-client}.

A second line of defense is cryptographic. Secure
aggregation~\cite{bonawitz2017secagg} hides individual client updates
behind a sum, leaving only the aggregate visible. Combined with DP,
secure aggregation reduces the adversary's effective observation to
\emph{neighbourhood-aggregated} updates rather than per-client ones,
weakening \TADI's parameter channel without addressing the structural
and organisational channels. \Fulcrum's allocation is compatible with
secure aggregation as a complementary defense; we discuss the
implications and leave a full empirical evaluation under secure
aggregation to future work (\S\ref{sec:discussion}).

A separate line on \emph{personalised} differential privacy
permits per-individual privacy budgets driven by user preferences.
J{\o}rgensen et~al.~\cite{jorgensen2015personalized} introduce
personalised DP for sample-release queries; subsequent work extends
this to ML training. Our setting is distinct: the asymmetry that
motivates per-client noise is structural (topology and organisational
position) rather than preferential, and the allocation is derived
from a worst-case MI bound rather than from declared user
preferences. The two lines are complementary; an integrated
allocation combining structural leverage with per-user preferences
is a natural extension we do not pursue here.

What none of these mechanisms address is \emph{which} clients should
receive how much noise when the topology breaks symmetry. Existing
implementations apply uniform noise multipliers across clients, implicitly
assuming a symmetric (\eg star) topology. Our defense generalises this:
when the topology is symmetric, our allocation reduces to uniform; when
it is not, the allocation balances the per-client privacy bound across
the asymmetric leverage profile.

\subsection{Topology in Federated Learning}
\label{sec:related-topology}

Topology-aware FL has been studied primarily as a performance and
scalability question. The earliest cross-device benchmark suite, LEAF~\cite{leaf2018},
demonstrates that natural client-partitioned datasets exhibit substantial
heterogeneity but does not consider topology beyond the implicit
star aggregator. Hierarchical FL~\cite{liu2020hierarchical,abad2020icassp}
introduces edge aggregators between clients and the cloud to reduce
communication cost. Decentralised gossip-based
schemes~\cite{roy2019braintorrent,hegedus2021decentralized} eliminate
the central aggregator entirely; Beilharz et~al.\ propose a DAG-based
extension~\cite{beilharz2021dag}. Briggs et~al.\ cluster clients by
local-update similarity to address non-IID data, producing a hierarchical
training structure that incidentally encodes data
similarity~\cite{briggs2020hierarchical}. The recent ``Field Guide''
survey~\cite{wang2021fieldguide} documents topology choice as one of the
canonical decisions in FL system design but treats it as orthogonal to
privacy.

The privacy literature, in turn, has implicitly assumed away topology.
Threat models almost universally assume a star configuration with a
single semi-honest aggregator. The attacker's view, when explicitly
modelled, consists of per-client gradients without structural metadata.
This implicit assumption is reasonable for cross-device deployments
(Gboard, smartphone keyboards) where every client is interchangeable;
it breaks down in cross-silo healthcare~\cite{flamby2022,sheller2020medicine,dayan2021covid},
multi-organisation industry consortia~\cite{melloddy2024}, and any
deployment that uses hierarchical or decentralised aggregation for
operational reasons. Our work fills this gap by treating the topology
as a first-class component of both the threat model and the defense.

\paragraph{Summary of positioning.}
We are the first, to our knowledge, to (i)~define a per-client
distributional inference attack with explicit channel ablations
isolating parameter, structural, and organisational information;
(ii)~prove an additive per-client MI bound that separates a
controllable mechanism term from an uncontrollable prior-coupling
floor; and
(iii)~derive a closed-form noise-allocation policy that strictly
improves on uniform DP-SGD whenever the topology breaks leverage
symmetry, with three principled proxies for instantiating leverage
in deployed regimes.
